% Source: transpose_convolution.tex
\begin{notechapter}
\chapter{Transpose Convolution}\label{ch:transpose-convolution}

\section{Introduction}
Transpose convolution, also known as deconvolution or
fractionally-strided convolution, is used to upsample
feature maps in neural networks. It increases the spatial
dimensions of the input, making it useful for tasks like
image segmentation and super-resolution.

\section{Input and Kernel}
Let’s start by defining the input and the kernel.

\subsection{Input \( x \) (size \( 2 \times 2 \))}
The input \( x \) is a \( 2 \times 2 \) matrix:
\[
x = \begin{bmatrix}
a & b \\
c & d
\end{bmatrix}
\]

\subsection{Kernel \( W \) (size \( 3 \times 3 \))}
The kernel \( W \) is a \( 3 \times 3 \) matrix:
\[
W = \begin{bmatrix}
w_{11} & w_{12} & w_{13} \\
w_{21} & w_{22} & w_{23} \\
w_{31} & w_{32} & w_{33}
\end{bmatrix}
\]

\section{Transpose Convolution as a ``Scatter'' Operation}

We take a stride \( s = 2 \) and a padding \( p = 1 \), as in
\texttt{torch.nn.ConvTranspose2d(1, 1, kernel\_size=3, stride=2, padding=1)}.
Where an ordinary convolution \emph{gathers} a weighted sum of a window of the
input into one output value, a transpose convolution \emph{scatters} each input
value into a window of the output:
\begin{enumerate}
    \item each input element \( x_{ij} \) is multiplied by the whole kernel \( W \);
    \item the resulting \( 3 \times 3 \) block is added to the output at position
          \( (s\,i, s\,j) \) (0-indexed), i.e.\ the blocks are placed \( s = 2 \) pixels apart,
          and overlapping contributions are summed;
    \item finally \( p = 1 \) row/column is cropped from each border of the result.
\end{enumerate}

\subsection{Output Size Formula}
Before cropping, the output has size \( (I - 1)s + K \); after cropping \( p \) on
each side, the output size \( O \) is:
\[
O = (I - 1) \times s + K - 2p
\]
where:
\begin{itemize}
    \item \( I \): Input size
    \item \( s \): Stride
    \item \( K \): Kernel size
    \item \( p \): Padding
\end{itemize}
(Frameworks add an optional \texttt{output\_padding} term to disambiguate the output size when
\( s > 1 \).) For \( I = 2 \), \( s = 2 \), \( K = 3 \), and \( p = 1 \):
\[
O = (2 - 1) \times 2 + 3 - 2 \times 1 = 3
\]
Thus, \( y \) is a \( 3 \times 3 \) feature map.

\subsection{Step-by-Step Computation}
Before cropping, the \( 5 \times 5 \) intermediate result is the sum of four shifted, scaled copies of
the kernel (with \( a \) at offset \( (0,0) \), \( b \) at \( (0,2) \), \( c \) at \( (2,0) \), \( d \) at \( (2,2) \)):
\[
Y_{\text{full}} =
\begin{bmatrix}
a w_{11} & a w_{12} & a w_{13} + b w_{11} & b w_{12} & b w_{13} \\
a w_{21} & a w_{22} & a w_{23} + b w_{21} & b w_{22} & b w_{23} \\
a w_{31} + c w_{11} & a w_{32} + c w_{12} & a w_{33} + b w_{31} + c w_{13} + d w_{11} & b w_{32} + d w_{12} & b w_{33} + d w_{13} \\
c w_{21} & c w_{22} & c w_{23} + d w_{21} & d w_{22} & d w_{23} \\
c w_{31} & c w_{32} & c w_{33} + d w_{31} & d w_{32} & d w_{33}
\end{bmatrix}.
\]
The overlaps (the third row and the third column) are where neighbouring input pixels both contribute.

\section{Full Output Matrix}
Cropping one row/column on each side gives the \( 3 \times 3 \) output:
\[
y = \begin{bmatrix}
w_{22}a & w_{23}a + w_{21}b & w_{22}b \\
w_{32}a + w_{12}c & w_{33}a + w_{31}b + w_{13}c + w_{11}d & w_{32}b + w_{12}d \\
w_{22}c & w_{23}c + w_{21}d & w_{22}d
\end{bmatrix}
\]
(This result was checked numerically against the equivalent formulation below.)

\section{Equivalent Formulation: Zero Insertion and Ordinary Convolution}
The same output is obtained with an ordinary (stride-1) convolution:
\begin{enumerate}
    \item insert \( s - 1 = 1 \) zero between neighbouring input elements:
    \[
    \begin{bmatrix} a & 0 & b \\ 0 & 0 & 0 \\ c & 0 & d \end{bmatrix};
    \]
    \item pad with \( K - 1 - p = 1 \) zero on each side (giving a \( 5 \times 5 \) matrix);
    \item slide the \emph{flipped} kernel (rotated by 180\(^\circ\)) with stride 1, computing sums of
          element-wise products (a cross-correlation), which gives a \( 3 \times 3 \) output.
\end{enumerate}
The zeros inserted between input values are the reason for the name
``fractionally-strided convolution''. Note that a transpose convolution is \emph{not} the inverse of a
convolution: it has the shape of the transpose of the convolution matrix (it maps a
\( 2 \times 2 \) map back to a \( 3 \times 3 \) map when the corresponding convolution maps \( 3 \times 3 \) to
\( 2 \times 2 \)), but not its values. This is also how the gradient of a strided convolution with respect to its
input is computed during backpropagation.

\section{Conclusion}
Transpose convolution is a powerful tool for upsampling.
By scattering each input value through a learnable kernel onto a finer grid
(equivalently, by inserting zeros and applying an ordinary convolution),
it produces larger feature maps, enabling tasks like
image generation and segmentation. When the stride does not divide the kernel size,
the uneven overlaps can produce ``checkerboard'' artifacts.
\end{notechapter}
